% ==================================================================
% id:       app-002
% topic:    שימושי הנגזרת
% subtopic: משפטי ערך ממוצע (רול, לגרנז׳, קושי)
% solution: full
% ==================================================================
תהא $f(x)$ פונקציה גזירה פעמיים בקטע \([-1,1]\) ונניח שמתקיים \\
\(f(-1) + f(1) = 2f(0)\). הוכיחו כי קיים \(x_0\) שעבורו \(f''\lr{x_0} = 0\),\\
כאשר \(x_0 \in (-1,1)\). \textbf{רמז:} משפטי רול ולגרנז'.
\hint{כתבו את הנתון כשוויון בין שתי מנות הפרשים, על שני חצאי הקטע. מה נותן משפט לגרנז' על כל חצי, ועל איזו פונקציה כדאי להפעיל אחר כך את משפט רול?}
\begin{solbox}
נציג את הנתון
$f(-1)+f(1)=2f(0)$
בצורה שתקל על פרושו:
$$.f(1)-f(0)=f(0)-f(-1)$$
כל אחד מהאגפים מזכיר ביטוי שמופיע במשפט לגרנז', אכן נציג
$$.\frac{f(1)-f(0)}{1}=\frac{f(0)-f(-1)}{1}$$
ממשפט לגרנז' ה"מופעל" על הפונקציה
$f$
בקטע
$[0,1]$,
קיים
$c_1\in(0,1)$
כך ש-
$$,f'(c_1)=\frac{f(1)-f(0)}{1}$$
כמוכן,
ממשפט לגרנז' ה"מופעל" על הפונקציה
$f$
בקטע
$[-1,0]$,
קיים
$c_2\in(-1,0)$
כך ש-
$$.f'(c_2)=\frac{f(0)-f(-1)}{1}$$
אנו רואים כי
$$.f'(c_1)=f'(c_2)$$
ממשפט רול המופעל על הפונקציה
$f'$
בקטע
$[c_2,c_1]\subset [-1,1]$,
קיים
$d\in(c_2,c_1)$
כך ש-
$$,f''(d)=0$$
כנדרש.
\end{solbox}
